% P2 / JCAA #280 -- Response to Reviewers, round 2 (revision v0.3).
% Drafted by Claude at the PI's request (2026-11-02, "Claude drafts, I approve"); the PI approves before upload.
% The reviewers' points are summarised, not quoted (public repository). Every statement was traced against
% submission_jcaa_v0.3.tex and checked by an independent reader on 2026-11-02 (findings applied the same day).
% RE-TRACE against the frozen manuscript after any later change.
\documentclass[11pt,a4paper]{article}
\usepackage[T1]{fontenc}
\usepackage[utf8]{inputenc}
\usepackage{lmodern}
\usepackage[margin=2.4cm]{geometry}
\usepackage{parskip}
\usepackage{enumitem}
\usepackage[hidelinks]{hyperref}
\setlist{itemsep=2pt, topsep=3pt}
\newcommand{\comment}[1]{\par\noindent\textit{Comment (summarised).} #1\par}
\newcommand{\response}{\par\noindent\textbf{Response.} }
\newcommand{\where}[1]{\par\noindent\textit{Where:} #1\par\medskip}

\begin{document}

\begin{center}
{\large\bfseries Response to the Reviewers}\\[4pt]
Journal of Computer Applications in Archaeology, submission \#280, second revision\\
\textit{An Evaluation Artefact in Presence--Background Archaeological Modelling: Evidence from East Java and a Corrected Comparison Protocol}
\end{center}

Dear Dr González-Pérez, dear Reviewers,

Thank you for the second round of comments. We have revised the text so that it no longer refers to the review
process or to an earlier version of the paper; we have explained the statistical terms for readers who are not
statisticians; we have added a new analysis to answer Reviewer~2's question about elevation and volcanic setting; and we
have redrawn Figures~1 and 5 and corrected the reference list.

While checking the data for this revision we also found a problem that neither reviewer raised: about a third of the
378 presence records used throughout the paper are modern monuments or misplaced records. We re-ran the real-data
analyses on a cleaned subset, and the result changed part of the paper's central claim. We describe this first, under
\emph{Changes not requested by the reviewers}, because it is the largest change in this version and the reviewers should
see it before the point-by-point answers.

Section, table and figure numbers refer to the revised manuscript.

% ------------------------------------------------------------------------------------------------
\section*{Changes not requested by the reviewers}

\subsection*{1. The presence records, a cleaned subset, and a restated central claim}

\textbf{What we found.} The 378 presence records come from OpenStreetMap (279), lists in the Indonesian and English
Wikipedia (85) and Wikidata (14). They sit on 335 distinct coordinates. About a third are not archaeological sites:
124 are modern monuments, statues or memorials, and 13 are misplaced or are not sites at all (for example, Central
Javanese temples carrying East Javanese coordinates). The previous version presented them, without qualification, as a
subset of East Java's archaeological record.

\textbf{What we did.} We defined a cleaning rule and re-ran the analyses on what it keeps. The rule removes 217 of the
378 records: the 85 rows whose coordinates were geocoded from their names (all rows taken from the Wikipedia lists),
the 124 modern monuments, and 8 temples of other regions, tourist objects and modern institutions; keeping one row per
coordinate removes nothing further. 161 sites remain (155 from OpenStreetMap, 6 from Wikidata). The design of this
re-analysis (E236) was committed to the public repository before it was run. It repeated the real-data contrasts the
paper's argument rests on (the common-background benchmark, the own- versus common-background matrix, the hard-negative
sweep, the priority maps and their seed agreement, and the elevation control), with the analysis code unchanged except
for the file of presences it reads.

\textbf{What changed in the results.}
\begin{itemize}
  \item On the 378 compiled records, redesigning the background with a common evaluation background held fixed
        contributes $-0.014$ AUC (interval $-0.021$ to $-0.005$ over five seeds), as before.
  \item On the 161 cleaned sites, the same redesign contributes $+0.051$ AUC. The pre-registered run used five seeds
        ($+0.0507$, positive in 14 of 15 seed$\times$algorithm cells). After that run we extended it to ten seeds, to
        meet our own rule of at least seven seeds for any reported value, and we report both: $+0.051$ ($+0.039$ to
        $+0.063$), positive in 29 of 30 cells, and two halves of five seeds that share no seed give $+0.0505$ and
        $+0.0516$. The interval is over seeds, not over sites; with 161 sites the detection floor is larger than for
        the 378 records, which the Limitations state.
  \item On the cleaned sites the hybrid design also beats the random design on every evaluation background (3/3
        algorithms on each of the four backgrounds), whereas on the compiled records it led only on its own background.
  \item What holds in both sets is the inflation from scoring a design on its own background: $+0.037$ on the compiled
        records (60/60 cells) and $+0.041$ on the cleaned sites (54/60). On the cleaned sites, own-background
        evaluation overstates the redesign gain by $+0.047$, close to the size of the gain itself.
  \item The selection criterion still has no interior optimum on the cleaned sites, and the target-group steps of the
        original sequence are not reproduced in either set. The elevation control gives the same primary result on
        both sets (below).
\end{itemize}

\textbf{How the claim was restated.} The previous version concluded that the rise of the original sequence was an
evaluation artefact and that, on a common evaluation background, redesigning the background adds nothing ($-0.014$
AUC). That statement now holds only for the compiled records. The paper now reports both sets and claims what holds in
both: own-background evaluation overstates the gain of a background design by about $+0.04$ to $+0.05$ AUC, and only a
common evaluation background shows how much real gain remains. The original sequence of experiments (E007--E013) is
kept on the 378 records, because it is the object under examination.

\textit{Where:} Abstract; Section~1.3; Section~2.1 (the records and the cleaning rule); new Section~3.3 and
Supplementary Table~S9 (both sets side by side); Sections~3.2 and 4.1; Conclusions; two new items and one amended
item in the Limitations (Section~4.6); Figure~1 (cleaned sites and other records drawn differently); Data Availability
(sources, licences and the cleaning script). The abstract therefore differs from the replacement abstract we sent to the
editor on 2 October.

\subsection*{2. Other corrections we made ourselves}
\begin{itemize}
  \item \textbf{Seed stability.} The previous version measured map agreement against a ten-seed ensemble that contains
        the seeds being averaged, which flatters the result. We now also report the agreement between ensembles that
        share no seed (Jaccard 0.76--0.86 for five seeds each on the compiled records, 0.70--0.77 on the cleaned
        sites) and present ``at least seven seeds'' as a floor, not as a guarantee of stability
        (Section~3.8, Figure~6, Table~1).
  \item \textbf{The target-group steps.} The previous version called the whole part of the rise attributed to
        background design an evaluation artefact. Only the hybrid step is shown to be one; the earlier target-group
        steps are not reproduced in our benchmark. The text now says so, and the heading of Section~3.2 now reads
        ``The ladder does not survive a common evaluation background'' (Sections~1.3 and 3.2).
  \item \textbf{``Pre-registered''.} Of the experiments E217--E224, only E224 had its design committed before the run;
        E218 and E220--E223 were run from written designs that entered the repository together with their results, and
        E217 and E219 had no written design. Section~2.6 states this, and ``pre-registered'' is now used only for E224
        and E235 (Sections~2.6 and 3.5).
  \item \textbf{Synthetic inflation by design.} The count ``343/360'' pooled the three designs. It is now given per
        design as well: hybrid 240/240 runs, random 52/60, target-group 51/60 (Section~3.2, Figure~3).
  \item \textbf{Robust core.} Its definition was stated wrongly; it is the set of cells in the top decile under all
        three background designs, and the caveat now gives the right reason why the density contrast is not a
        validation (the same sites trained the models) (Section~3.8).
  \item \textbf{Share of hard negatives.} The realised share of hard negatives in the hybrid design (0.62, against a
        target of 0.30) was attributed to the target-group pool. The proportion of such points in the target-group
        background (0.485) is close to that of the random background (0.503), so the explanation is now the
        composition of the whole background pool (Section~2.3).
  \item \textbf{Test 3.} The previous version used a performance check across road-distance quartiles (``Test 3'').
        It was scored on each design's own background, so it is subject to the artefact this paper describes; the
        text now says so and does not use it (Section~2.4).
  \item \textbf{Data description.} The Data Availability statement listed a source that contributed no records and
        gave a wrong licence for Wikidata; both are corrected.
  \item \textbf{Authors' contributions and Declaration of AI use.} The Authors' Contributions now say that the
        analyses were implemented with AI assistance, and the Declaration of AI Use states how the analyses and the
        text were produced, in the terms of our message to the editor of 2 November.
\end{itemize}

% ------------------------------------------------------------------------------------------------
\section*{Reviewer 1}

\textbf{R1-1}
\comment{The paper should be self-contained: no mention of reviewers, requests, responses, revisions or an earlier
version; the story of the rework could go in the acknowledgements.}
\response We removed or rewrote every passage of the main text that referred to the reviewers, to the review process or
to an earlier version of the paper (38 in the previous version), including the abstract. The development of the
analysis is now presented as the study's own sequence of experiments (Sections~1.3, 2.3 and 3.1), and the later
experiments as tests of that sequence. Outside the main text, the Acknowledgements keep the conventional thanks to the
editor and reviewers, and the Declaration of AI Use describes how the earlier drafts were produced, as the journal's AI
policy requires.
\where{whole text; Abstract; Acknowledgements; Declaration of AI Use.}

\textbf{R1-2}
\comment{The research question should be stated earlier, more clearly and more pedagogically for readers who are not
statisticians.}
\response The first paragraph of the Introduction now states the question in plain terms: a change in how background
points are selected can raise the cross-validated AUC without the model becoming better, because the change can make the
evaluation task easier; the study tests which of the two is happening.
\where{Section~1, first paragraph.}

\textbf{R1-3}
\comment{The statistical methods need more explanation for non-specialists.}
\response We added Table~1, \emph{Key terms}, which defines in plain language the terms the paper relies on:
presence--background model, background design, hard negatives, evaluation background, availability domain, AUC, TSS,
continuous Boyce index, spatial block cross-validation, seed ensemble, Jaccard overlap and coarsened exact matching.
\where{Table~1 (Section~1.1).}

\textbf{R1-4}
\comment{Three expressions are hard to understand.}
\response All three are replaced by literal wording: ``we sweep the hard-negative dial'' now reads ``we vary the share
of hard negatives in the background from 0 to 1'' (Section~2.6); ``how large is judging at home?'' now reads ``how large
is own-background inflation?'' (Section~3.2); and ``map Jaccard said no'' now reads ``the pre-registered primary metric
was map Jaccard, and it did not support that reading'' (Section~3.6). We also removed the rest of the same family of
expressions throughout the text and in the figures (for example ``dial'', ``home-court'' and ``judged at home''). The
word ``ladder'' is kept for the original sequence of seven experiments and is defined where it first appears
(Section~1.3).
\where{Sections~1.3, 2.6, 3.2 and 3.6; Figures~2--4.}

\textbf{R1-5}
\comment{The section sign is unusual.}
\response All section references now read ``Section~n''.
\where{whole text.}

\textbf{R1-6}
\comment{The bibliography is inconsistently formatted and some names are misspelt; the Comer et al.\ reference is
wrong.}
\response We checked every reference that has a DOI against Crossref. Three entries carried wrong metadata and are
corrected: Comer et al.\ (2023) (authors, title and volume), Wang, Shi \& Oguchi (2023) (authors and title) and
Castiello \& Tonini (2021) (title and the first author's given name). The sentence that cited these three papers
together did not describe them accurately; it has been rewritten from what the papers report (Section~1.3). Name
spellings were corrected, DOIs were added where they exist, title capitalisation was made consistent, and the whole list
now follows the journal's Harvard style.
\where{References; Section~1.3.}

\textbf{R1-7}
\comment{Figure~1 should show the study area; it is a scatter plot whose axis values cannot be read.}
\response Figure~1 is now a relief map drawn from the elevation model used by the models, with the coastline, a
latitude--longitude grid, a scale bar, a north arrow and a locator map of Indonesia. It shows the presence records (the
161 cleaned sites and the other 217 records drawn differently), the 13 volcanic centres inside the study area, and the
20~km and 40~km distance lines used in the new elevation control.
\where{Figure~1.}

\textbf{R1-8}
\comment{The legend of Figure~5 should not sit on the map.}
\response The legend is now below the map. The map is drawn as a raster with the same coastline, grid and scale bar as
Figure~1.
\where{Figure~5.}

\textbf{R1-9}
\comment{The text should be cleaned of the history of the previous version.}
\response As for R1-1.

% ------------------------------------------------------------------------------------------------
\section*{Reviewer 2}

\textbf{R2-1}
\comment{It is still unclear whether the model responds to volcanic areas or simply to high elevation; comparing
volcanic land with non-volcanic land at the same elevation would settle it.}
\response We ran the comparison you suggest as a new analysis (E235), with its design committed to the public
repository before it was run. Land is ``near'' if it lies within 20~km of a Quaternary volcanic centre and ``far'' if it
lies 40~km or more from one; within each 200~m elevation band the two are matched on slope, terrain ruggedness and
topographic wetness, and uncertainty comes from a spatial block bootstrap (1,000 replicates). The result has two parts.
\begin{itemize}
  \item \textbf{Above about 1,000~m the comparison cannot be made in East Java.} Of the roughly 3,435~km$^2$ of land
        above 1,000~m, 97.5\% lies within 20~km of a Quaternary volcanic centre and only 0.4\% lies 40~km or more
        from one (about 15~km$^2$ at 1,000--1,200~m, none above 1,200~m). Elevation and volcanic setting therefore
        cannot be separated there, and the paper now says so and attributes the pattern above 1,000~m to neither.
  \item \textbf{Below 1,000~m no direction is claimed.} With XGBoost, the primary learner, matched suitability is
        $+0.049$ higher near volcanic centres for the random design, $+0.049$ for the target-group design and $+0.060$
        for the hybrid design. Under the pre-registered rule the hybrid result alone meets the criterion for a
        volcano-specific difference, and the summary across the three designs is inconclusive. For the random design
        the two pre-registered sensitivity definitions give $-0.002$ and $+0.016$, and post hoc variants give $+0.005$
        to $+0.07$. The far uplands lie mostly in the Southern Mountains, an older (Eocene--Miocene) volcanic arc, so
        they are not strictly non-volcanic land, and the text states this.
\end{itemize}
No model input carries volcanic information. We removed the earlier terrain-matched comparison, which had only two sites
in its non-volcanic arm, and the wording that its result was consistent with a volcano-specific effect; the new analysis
replaces it, and the paper no longer implies that the model's response is volcano-specific. On the cleaned sites the
primary result is the same (the hybrid design alone meets the criterion; the summary is inconclusive), while the two
sensitivity definitions exchange their labels between inconclusive and not specific (Section~3.3).
\where{new Section~3.10; Figure~7; Supplementary Tables~S7--S8; Section~4.6 (Limitations); Section~3.3.}

\textbf{R2-2}
\comment{Without topography in Figure~1, readers cannot see how the model relates to elevation.}
\response The new Figure~1 shows elevation (a hillshaded elevation model with an elevation scale) together with the
distance lines used in the new control.
\where{Figure~1.}

\bigskip
We thank both reviewers again. The remaining errors are ours.

\medskip
Mukhlis Amien and Go Frendi Gunawan\\
Universitas Bhinneka Nusantara

\end{document}
